\documentclass[11pt]{article}

\usepackage[margin=1in]{geometry}
\usepackage{amsmath,amssymb,amsthm,mathtools}
\usepackage[hidelinks]{hyperref}

\newtheorem{theorem}{Theorem}[section]
\newtheorem{lemma}[theorem]{Lemma}
\newtheorem{proposition}[theorem]{Proposition}
\newtheorem{corollary}[theorem]{Corollary}
\theoremstyle{remark}
\newtheorem{remark}[theorem]{Remark}

\newcommand{\splus}{s^+}
\newcommand{\sminus}{s^-}
\newcommand{\sig}{\sigma}
\newcommand{\Arg}{\operatorname{Arg}}

\title{Cacti with Arbitrarily Many Triangles and One Hostile Cycle}
\author{}
\date{28 July 2026}

\begin{document}
\maketitle

\begin{abstract}
Let $G$ be a connected cactus whose cyclic blocks are $r\geq1$ triangles and
one cycle $C_q$, where $q=4k+1\geq5$.  The cyclic blocks may share cut
vertices in arbitrary incidence topology, bridge paths may have arbitrary
length, and arbitrary finite trees may be attached throughout.  We prove
\[
 \splus(G)>|V(G)|.
\]
The proof is rank-uniform.  A maximum-packing Voronoi partition proves strict
positivity for every triangular cactus and yields a rooted guard theorem for
one hostile cycle.  At the global level, actual bridges remove triangular leaf
clusters inductively.  In the remaining shared-cut cluster, one interface is
handled by the rooted theorem, while two or more interfaces allow the hostile
cycle itself to be split into consecutive intervals, leaving only strict
triangular induced territories.  Every cut, connector remnant, and attached
tree receives exactly one owner.
\end{abstract}

\section{Statement and basic facts}

All graphs are finite, simple, and undirected.  For the adjacency eigenvalues
$\lambda_1,\ldots,\lambda_n$ of a graph $H$, put
\[
 \splus(H)=\sum_{\lambda_i>0}\lambda_i^2,
 \qquad
 \sminus(H)=\sum_{\lambda_i<0}\lambda_i^2,
 \qquad
 \sig(H)=\splus(H)-|V(H)|.
\]
A cactus is a graph whose blocks are bridges, isolated vertices, or cycles.
Its cyclomatic number is the number of cyclic blocks.

\begin{theorem}\label{thm:main}
Let $G$ be a connected cactus whose cyclic blocks are $r\geq1$ copies of
$C_3$ and one copy of $Q=C_q$, where $q=4k+1\geq5$.  Allow arbitrary bridge
connectors and arbitrary finite trees attached at arbitrary vertices.  Then
\[
 \boxed{\splus(G)>|V(G)|}.
\]
\end{theorem}

We use square-energy superadditivity on induced vertex partitions.

\begin{lemma}[Induced-partition superadditivity]\label{lem:super}
If $V(H)=V_1\dot\cup\cdots\dot\cup V_t$, then
\[
 \splus(H)\geq\sum_{j=1}^t\splus(H[V_j]),
 \qquad
 \sig(H)\geq\sum_{j=1}^t\sig(H[V_j]).
\]
\end{lemma}

\begin{proof}
For a real symmetric matrix $A$, its positive spectral part satisfies
\[
 \operatorname{tr}(A_+^2)
 =\max_{X\succeq0}\bigl(2\operatorname{tr}(AX)-\operatorname{tr}(X^2)\bigr).
\]
This follows by diagonalizing $A$: deleting the off-diagonal entries of the
positive semidefinite test matrix preserves positive semidefiniteness and only
increases the objective, after which the scalar maximizers are
$x_i=\max(\lambda_i,0)$.

Order the vertices by the partition, let $A=A(H)$, and put
$B=\operatorname{diag}(A(H[V_1]),\ldots,A(H[V_t]))$.  Testing the variational
formula with $X=B_+$ gives
\[
 \splus(H)\geq2\operatorname{tr}(AB_+)-\operatorname{tr}(B_+^2)
 =\operatorname{tr}(B_+^2)=\sum_j\splus(H[V_j]),
\]
because $B_+$ is block diagonal and the off-block entries of $A$ do not
contribute to the trace.  Subtracting the partitioned vertex count proves the
second inequality.
\end{proof}

\section{Triangular territories}

For a graph $F$ of order $n$, put
\[
 \Psi_F(t)=i^{-n}\det(itI-A(F))
 =\prod_{j=1}^n(t+i\lambda_j),\qquad t>0,
\]
and define its continuous phase, tending to zero at infinity, by
\[
 \Theta_F(t)=\sum_{j=1}^n\arctan\frac{\lambda_j}{t}.
\]

\begin{lemma}[Signed Coulson identity]\label{lem:coulson}
For every graph $F$,
\[
 \splus(F)-\sminus(F)
 =-\frac4\pi\int_0^\infty t\Theta_F(t)\,dt.
\]
\end{lemma}

\begin{proof}
For real $\lambda$,
\[
 \lambda|\lambda|=\frac2\pi\int_0^\infty
       \frac{\lambda^3}{\lambda^2+t^2}\,dt.
\]
Also $\Theta_F'(t)=-\sum_j\lambda_j/(\lambda_j^2+t^2)$.  Since
$\sum_j\lambda_j=\operatorname{tr}A(F)=0$,
\[
 \sum_j\frac{\lambda_j^3}{\lambda_j^2+t^2}=t^2\Theta_F'(t).
\]
Sum the scalar identity and integrate by parts.  The boundary term vanishes:
$t^2\Theta_F(t)\to0$ at zero, while the arctangent expansion gives
$\Theta_F(t)=O(t^{-3})$ at infinity.  These estimates also give convergence.
\end{proof}

The cycle-packing number $\nu(H)$ is the maximum number of pairwise
vertex-disjoint cycles in $H$.

\begin{lemma}[Packing-one triangular packet]\label{lem:packing-one}
Let $H$ be a connected cactus whose $b\geq1$ cyclic blocks are triangles and
whose cycle-packing number is one.  Arbitrary trees may be attached.  Then
\[
 \splus(H)>\sminus(H),\qquad \sig(H)>b-1\geq0.
\]
\end{lemma}

\begin{proof}
For $t>0$, let
\[
 Z_F(t)=\sum_M t^{|V(F)|-2|M|},
 \qquad
 \Psi_H(t)=i^{-|V(H)|}\det(itI-A(H)),
\]
where $M$ ranges over matchings.  Grouping the Sachs expansion by collections
of vertex-disjoint cycles, a triangle has multiplier $-2i$.  Because no two
triangles are disjoint,
\[
 \Psi_H(t)=Z_H(t)-2i\sum_T Z_{H-V(T)}(t).
\]
Every displayed matching partition is positive.  Hence the continuous
argument $\Theta_H(t)$ of $\Psi_H(t)$, tending to zero at infinity, lies
strictly in $(-\pi/2,0)$.  Lemma~\ref{lem:coulson}
\[
 \splus(H)-\sminus(H)
 =-\frac4\pi\int_0^\infty t\Theta_H(t)\,dt
\]
therefore gives strict positivity.  Since
$|E(H)|=|V(H)|-1+b$ and
$\splus(H)+\sminus(H)=2|E(H)|$, the surplus exceeds $b-1$.
\end{proof}

\begin{lemma}[Maximum-packing territories]\label{lem:voronoi}
Let $A$ be a connected cactus all of whose cyclic blocks are triangles.  Let
$T_0,\ldots,T_{m-1}$ be a maximum-cardinality collection of pairwise
vertex-disjoint triangles.  Assign each vertex $v$ to the lexicographically
least pair
\[
 \bigl(d_A(v,V(T_j)),j\bigr),
\]
and let $A_j$ be the induced graph on the vertices assigned to $j$.  Then each
$A_j$ is connected, contains $T_j$, and has cycle-packing number one.
If a root $z$ is prescribed and $T_0$ is chosen nearest $z$, then $z\in A_0$.
\end{lemma}

\begin{proof}
A predecessor of an assigned vertex on a shortest path to its selected cycle
has the same owner: distances are one-Lipschitz, and the priority rule resolves
ties consistently.  Thus every territory is connected and contains its
selected cycle.  If $A_j$ contained two disjoint cycles $D,E$, then
\[
 D,E,\{T_h:h\ne j\}
\]
would be a packing larger by one.  Notice that $D,E$ replace $T_j$; they need
not also be disjoint from it.  The root assertion follows from choosing $T_0$
with minimum root distance and giving it first priority.
\end{proof}

\begin{corollary}[All triangular cacti]\label{cor:triangular}
Every connected cactus containing at least one cyclic block, all cyclic blocks
being triangles, satisfies $\sig(A)>0$, with arbitrary attached trees.
\end{corollary}

\begin{proof}
Apply Lemma~\ref{lem:voronoi}.  Every territory is a packing-one triangular
cactus retaining at least one triangle, so Lemma~\ref{lem:packing-one} gives
strictly positive surplus to every territory.  Sum by
Lemma~\ref{lem:super}.
\end{proof}

\section{One rooted hostile cycle}

Put
\[
 \delta_q=\sec\frac{\pi}{q}-1.
\]
For $q\geq5$, $0<\delta_q<1$.

\begin{lemma}[Packing-one hostile packet]\label{lem:hostile-packet}
Let $H$ be a connected cactus whose cyclic blocks are $Q=C_q$ and $a\geq1$
triangles, where $q=1\pmod4$.  Suppose no two triangles are vertex-disjoint.
The cycle $Q$ may meet the triangular part at one cut vertex or be joined to a
distinguished root of that part by one path.  Arbitrary trees may be attached
anywhere.  Then
\[
 \splus(H)-\sminus(H)>-2\delta_q,
 \qquad
 \sig(H)>a-\delta_q\geq1-\delta_q>0.
\]
\end{lemma}

\begin{proof}
Let the \emph{spine} be the subgraph consisting of all cyclic blocks and the
unique minimal paths joining them.  Every component outside it is a tree with
one spine attachment; two attachments would create another cycle.  Orient each
off-spine branch toward the spine.  If $u\to v$ is an oriented branch edge,
let $T_{u\to v}$ be the component below $u$, including $u$, and define
\[
 M_{u\to v}(t)=\frac{Z_{T_{u\to v}}(t)}
                       {Z_{T_{u\to v}-u}(t)}.
\]
Splitting matchings according to whether $u$ is unmatched or matched to one
child proves the exact recursion
\[
 M_{u\to v}(t)=t+\sum_{w\text{ child of }u}
                   \frac1{M_{w\to u}(t)}\geq t.                 \tag{1}
\]
The leaf case is $M=t$.  Extracting all denominators recursively gives a
common positive factor $K(t)$ and replaces the unmatched-vertex activity at a
spine vertex by
\[
 \alpha_v(t)=t+\sum_{u\text{ branch child of }v}
                    \frac1{M_{u\to v}(t)}=t+y_v(t),
 \qquad y_v(t)\geq0.                                           \tag{2}
\]
This is an exact matching decomposition.

For positive activities $\alpha=(\alpha_v)$, define
\[
 Z_F(\alpha)=\sum_{M\text{ matching of }F}
                  \prod_{v\text{ unmatched by }M}\alpha_v.
\]
It is positive and coordinatewise nondecreasing.  Tree elimination says that
each Sachs matching carrier is $K(t)$ times the corresponding weighted spine
carrier.  Let $S=V(Q)$.
The grouped Sachs expansion has the form
\[
 \frac{\Psi_H(t)}{K(t)}=R+2i(B-A),
\]
where
\[
 B=Z_{\mathrm{spine}-S}(\alpha)>0,
 \quad
 A=\sum_T Z_{\mathrm{spine}-V(T)}(\alpha)>0,
\]
and
\[
 R=Z_{\mathrm{spine}}(\alpha)
   +4\sum_{T\cap Q=\varnothing}
       Z_{\mathrm{spine}-(S\cup V(T))}(\alpha)>0.
\]
The normalized Sachs multiplier of $C_\ell$ is $-2i^{-\ell}$.  Thus a triangle
contributes $-2i$, while $q=1\pmod4$ makes $Q$ contribute $+2i$.  A Sachs
collection uses vertex-disjoint cycles, so it contains at most one triangle.
The empty collection gives $Z_{\mathrm{spine}}(\alpha)$, singleton $Q$ and
triangle collections give $2iB$ and $-2iA$, and a disjoint $Q$--triangle pair
gives $(2i)(-2i)=4$.  These exhaust the expansion and prove the displayed
formulas.

Let $Z_q(t)$ denote the unweighted matching partition of $C_q$.  Partition
spine matchings according to whether they use an edge between $S$ and its
complement.  Those using no such edge factor into a matching on $Q$ and one on
the vertex-deleted remainder; all remaining matchings have nonnegative weight.
Thus
\[
 Z_{\mathrm{spine}}(\alpha)=Z_Q(\alpha|S)B+E,
 \qquad E\geq0.
\]
Coefficient positivity and $\alpha_v\geq t$ give
$Z_Q(\alpha|S)=Z_q(t)+L$ with $L\geq0$.  Therefore
\[
 R-Z_q(t)(B-A)
 =E+LB+Z_q(t)A
 +4\sum_{T\cap Q=\varnothing}
   Z_{\mathrm{spine}-(S\cup V(T))}(\alpha)>0.             \tag{1}
\]
Since $K(t)>0$ it changes no phase.  Since $R>0$, the continuous phase is the
principal value
\[
 \Theta_H(t)=\arctan\frac{2(B-A)}R.
\]
The isolated hostile cycle has
$\Psi_Q(t)=Z_q(t)+2i$ and phase
$\theta_q(t)=\arctan(2/Z_q(t))$.  If $B-A\leq0$, then
$\Theta_H\leq0<\theta_q$; if $B-A>0$, inequality (1) gives the same strict
comparison.  Thus $\Theta_H(t)<\theta_q(t)$ for all $t>0$.

Lemma~\ref{lem:coulson} and the exact cycle value
\[
 \splus(C_q)-\sminus(C_q)=-2\left(\sec\frac\pi q-1\right)
\]
give the first displayed inequality.  Finally
$|E(H)|=|V(H)|+a$ and
$\splus(H)+\sminus(H)=2|E(H)|$, proving the surplus bound.
\end{proof}

\begin{lemma}[Exact hostile-cycle value]\label{lem:qvalue}
For $q=4k+1\geq5$,
\[
 \splus(C_q)-\sminus(C_q)
 =-2\left(\sec\frac\pi q-1\right).
\]
\end{lemma}

\begin{proof}
The eigenvalues are $2\cos(2\pi j/q)$, $0\leq j<q$.  The positive ones have
$j=0,1,\ldots,k$ and $j=q-k,\ldots,q-1$, so
\[
 \splus(C_q)=4\left(1+2\sum_{j=1}^k
                   \cos^2\frac{2\pi j}{q}\right).
\]
Using $\cos^2x=(1+\cos2x)/2$ and
\[
 \sum_{j=1}^k\cos(jx)
 =\frac{\sin(kx/2)\cos((k+1)x/2)}{\sin(x/2)}
\]
at $x=4\pi/q$ gives
$\splus(C_q)=q+1-\sec(\pi/q)$.  Since
$\splus(C_q)+\sminus(C_q)=2q$, the claimed difference follows.
\end{proof}

\begin{theorem}[Rooted all-rank guard]\label{thm:rooted}
Let $A$ be an arbitrary connected triangular cactus with at least one triangle,
rooted at $z$.  Attach $Q=C_q$, $q=1\pmod4$, to $z$ through exactly one
interface: identify a vertex of $Q$ with $z$, or join them by one path.
Allow arbitrary attached trees.  Then
\[
 \sig(G)>1-\delta_q>0.
\]
\end{theorem}

\begin{proof}
Apply Lemma~\ref{lem:voronoi} to a maximum triangle packing of $A$, choosing
the root territory first.  Add $Q$, the whole connector, and all trees rooted
there to the root territory.  Its retained triangles have packing one and
include the selected root triangle, so Lemma~\ref{lem:hostile-packet} gives
surplus $>1-\delta_q$.  Every other induced territory is a packing-one
triangular cactus and has positive surplus by Lemma~\ref{lem:packing-one}.
The unique-interface hypothesis ensures that no connector vertex must have two
owners.  Sum by Lemma~\ref{lem:super}.
\end{proof}

\section{Shared-cut clusters and ownership}

Two cyclic blocks are adjacent when they share a cut vertex.  A
\emph{shared-cut cluster} is a maximal connected component under this
adjacency.  In the cactus block--cut tree, the \emph{support} of a cluster is
the connected subtree consisting of its cyclic-block nodes and every
cut-vertex node joining adjacent cyclic blocks in the cluster.  Distinct
supports are disjoint, since a cut belonging to both would join the clusters.
Contract each full support to one cluster node.  In the resulting tree take
the minimal subtree spanning all cluster nodes and suppress each noncluster
vertex of degree two.  This is the \emph{reduced cluster tree}.  Every leaf is
a cluster node, and every reduced edge contains a nonempty chain of actual
bridge blocks; otherwise its endpoint cycles would lie in one cluster.

When one actual bridge in such a chain is cut, its two endpoint components give
a vertex partition into connected induced cacti.  Every remaining connector
segment and every off-hull tree follows the side containing its unique
attachment.  An off-hull component cannot have two hull attachments: together
with the hull path they would create another cycle.  This is the ownership rule
used below.

For one shared-cut cluster, form its bipartite cycle--cut incidence graph $I$:
cycle nodes are cyclic blocks, cut nodes are vertices belonging to at least two
cyclic blocks, and containment gives incidence.  This graph is a tree.  Indeed,
an incidence cycle would realize either a cycle using several blocks or two
blocks meeting in more than one vertex, both impossible in a cactus.

\begin{lemma}[Splitting the distinguished cycle]\label{lem:qsplit}
Suppose one shared-cut cluster consists of $Q$ and triangles.  If $Q$ is
incident in $I$ with $d\geq2$ distinct cut nodes, then the graph has an induced
vertex partition into $d$ connected nonempty triangular cacti.
\end{lemma}

\begin{proof}
Let $x_1,\ldots,x_d$ be the distinct occupied vertices of $Q$ in cyclic order.
Partition $V(Q)$ into $d$ nonempty proper consecutive intervals $J_i$, each
containing $x_i$ and no other occupied vertex.  Consecutive occupied vertices
may receive singleton intervals; the construction also works when all vertices
of $Q$ are occupied.

Deleting the $Q$ node from $I$ gives exactly $d$ components $K_i$, with $x_i$
in $K_i$.  The \emph{realization} of $K_i$ consists of all private vertices of
its cycle blocks and all its cut vertices.  Give one territory the interval
$J_i$, the complete realization of $K_i$, and every tree whose unique
attachment it owns.  Since $x_i$ is a
shared cyclic cut, $K_i$ contains at least one triangle.  Each territory is
connected through $x_i$.  The intervals partition $Q$, the components $K_i$
partition the remaining incidence tree, and all multiway branches at one cut
remain with one owner.  Thus the territories are induced, disjoint, and
exhaustive.  Every $J_i$ is a proper path, so $Q$ is retained nowhere; all
retained cycles are triangles.
\end{proof}

\section{Proof of the main theorem}

\begin{proof}[Proof of Theorem~\ref{thm:main}]
Induct on the number $N$ of shared-cut clusters.

Suppose first that $N>1$.  Only one cluster contains $Q$.  The reduced cluster
tree has at least two leaves, so it has a leaf cluster $L$ not containing $Q$.
Cut the first actual bridge away from $L$, using the ownership rule above.  One
induced side $A$ is a connected triangular cactus, hence
$\sig(A)>0$ by Corollary~\ref{cor:triangular}.  The other side $B$ contains
$Q$ and all remaining clusters.

If $B$ retains at least one triangle, it is a smaller instance of the theorem,
so induction and Lemma~\ref{lem:super} give strict positivity.  If $B$ retains
no triangle, there is no other triangular cluster, so $N=2$, $L$ contains all
triangles, and the original graph consists of a connected triangular cactus
joined to $Q$ through exactly one bridge
interface.  Theorem~\ref{thm:rooted} closes this endpoint directly.  This avoids
the invalid step of asking qualitative triangular strictness to pay the
negative surplus of an isolated hostile cycle.

It remains to treat $N=1$.  Let $d$ be the degree of the $Q$ node in the
cycle--cut incidence tree.  Since at least one triangle is present and the
cluster is connected, $d\ne0$.

If $d=1$, deleting the $Q$ node leaves one connected triangular cactus rooted
at the unique shared cut.  The cycles need not all contain that root, so a
common-cut bouquet theorem would be insufficient; Theorem~\ref{thm:rooted}
applies exactly and gives strict positivity.

If $d\geq2$, Lemma~\ref{lem:qsplit} partitions $G$ into connected nonempty
triangular cacti.  Every part has positive surplus by
Corollary~\ref{cor:triangular}, and Lemma~\ref{lem:super} completes the proof.
\end{proof}

\section{Scope}

The proof is uniform in the number of triangular blocks and does not enumerate
incidence trees.  It depends essentially on there being one distinguished
hostile cycle.  With two pentagons, a maximum-packing partition can retain the
two hostile cycles while splitting every triangular guard; already the
incidence path $C_5-x-C_3-y-C_5$ exhibits this ownership obstruction.  Thus the
present theorem does not claim the arbitrary $T^rC_5C_5$ family.

\begin{thebibliography}{9}
\bibitem{AKMPZ}
S.~Akbari, H.~Kumar, B.~Mohar, S.~Pragada, and X.~Zhang,
\emph{Refinement of a conjecture on positive square energy of graphs},
arXiv:2506.07264v1, 2025.

\bibitem{Coulson}
C.~A. Coulson,
\emph{On the calculation of the energy in unsaturated hydrocarbon molecules},
Proc. Cambridge Philos. Soc. 36 (1940), 201--203.
\end{thebibliography}

\end{document}
